\section{Experimental Setup}
\label{sec:experimental-setup}
\subsection{Datasets}
We run our experiments on two public bone X-ray datasets. BTXRD is a primary bone tumor dataset with polygon annotations supporting classification, localization, and segmentation \cite{b7}. FracAtlas is a fracture dataset with the same label types \cite{b8}; we convert its segmentation annotations to the same polygon format used for BTXRD before use, and the conversion is released with the reproducibility package. Both datasets are split at the image level, so that every polygon belonging to the same image always falls in the same split (Table~\ref{tab:splits}).

\begin{table}[t]
\caption{Image-level train / validation / test splits. Each lesion polygon is one promptable sample.}
\label{tab:splits}
\centering
\scriptsize
\resizebox{\columnwidth}{!}{%
\begin{tabular}{lcccccc}
\toprule
& \multicolumn{3}{c}{Images} & \multicolumn{3}{c}{Lesion polygons} \\
\cmidrule(lr){2-4}\cmidrule(lr){5-7}
Dataset & Train & Val & Test & Train & Val & Test \\
\midrule
BTXRD & 1{,}493 & 187 & 187 & 1{,}848 & 238 & 232 \\
FracAtlas & 573 & 72 & 72 & 730 & 95 & 92 \\
\bottomrule
\end{tabular}
}
\end{table}

Table~\ref{tab:splits} contains 917 converted lesion polygons for FracAtlas (730 + 95 + 92), whereas the original FracAtlas release reports 922 fracture instances. These totals have not been reconciled at the annotation-ID level. The reported segmentation results therefore refer to the converted annotation set summarized in Table~\ref{tab:splits}. Because the public metadata do not provide enough information to verify strict patient-level separation, images from the same patient could in principle still fall into different splits; we regard this as a limitation of this evaluation. In image-level evaluation, if an image has multiple lesions, the model is run separately for each box, and the resulting probability maps are merged by a pixelwise maximum before thresholding at 0.5; the ground-truth masks are likewise merged by union.

Input images are resized isotropically to preserve aspect ratio, then padded to a square resolution of $128\times128$, $256\times256$, or $512\times512$, depending on the configuration. Images, masks, and prompt maps all undergo this same geometric preprocessing.

\subsection{Training Details}
\label{sec:training}
All CNN models were implemented in PyTorch and trained on a single NVIDIA Tesla T4 GPU with 16 GB memory. PGA-UNet and Attention U-Net were trained separately for each dataset and each resolution; in other words, this article studies one shared architecture family instantiated as separate models, with separate parameter sets for BTXRD and FracAtlas, rather than training a single joint model on both. We used the AdamW optimizer with learning rate $10^{-4}$ and weight decay $10^{-4}$; the loss was the sum of binary cross-entropy and Dice loss. Because the segmentation targets are often small, we additionally tested training PGA-UNet at $512\times512$ with the Dice term replaced by a size-conditioned Tversky loss, and in a separate run, by a Focal Dice loss, with every other setting held fixed. These two replacements are mutually exclusive and are reported later only as a negative control on the training objective. Batch size was 4, training ran for at most 150 epochs, with early stopping at patience 15. All main and ablation experiments share the same fixed training seed, 22120196, which seeds Python, NumPy, and PyTorch on both CPU and CUDA. The primary model-selection criterion was image-level merged Dice, measured on the validation set under the \texttt{center\_shift} condition.

During training, we applied augmentation consisting of horizontal flipping and random rotation within $\pm15^\circ$. PGA-UNet is trained with \texttt{center\_mixed}: \texttt{center\_shift} is chosen with probability 0.8 and \texttt{center\_zoom} with probability 0.2. Because the prompts are all simulated from labeled lesion regions, this is a prompted-segmentation protocol under controlled conditions; the article should be read accordingly, namely as a study of simulated lesion-covering boxes around user-identified lesions, not as an evaluation of unconstrained lesion detection or of boxes fully drawn by a clinician in real practice.

For the fine-tuned SAM-Med2D comparison, the pretrained image encoder was frozen, with only its lightweight Adapter layers trained further, while the prompt encoder and mask decoder were fully updated on each dataset, under the same split protocol. SAM-Med2D was fine-tuned with box prompts only, using the same center-scaled covering and off-center protocol as PGA-UNet; the point-prompt branch and iterative point refinement were both disabled. This box construction differs from the small random-jitter \texttt{get\_boxes\_from\_mask} sampling used by the original SAM-Med2D authors, which we did not apply here. Both validation during SAM-Med2D fine-tuning and the final test-split evaluation report both prompt conditions separately; batch size was 4, at most 150 epochs, early stopping patience 30, and learning rate $10^{-5}$. This matched protocol reduces differences caused by prompt distribution and resolution, but cannot eliminate differences in model scale, initialization, or optimization.

We additionally ran stability experiments using Monte Carlo cross-validation, that is, repeated random image-level splits, as opposed to strict non-overlapping $k$-fold partitioning. These experiments keep the training seed fixed at 22120196 and vary only the split seed, across the values 1, 2, 3, and 4.

Two details are worth noting when interpreting the evaluation results. First, prompt construction itself does not depend on resolution: the plateau heatmap is built and smoothed with a fixed $31\times31$ Gaussian kernel at the original image resolution, and only then resized and padded down to $128\times128$, $256\times256$, or $512\times512$ together with the image. In other words, the same absolute blur is compressed differently at each target resolution, rather than being parameterized separately for each. Second, the Monte Carlo experiments are only auxiliary stability evidence, not intended to replace the fixed train/validation/test split used throughout the main results tables.

\subsection{Baselines and Metrics}
We compare PGA-UNet against an automatic baseline, Attention U-Net \cite{b4}, and three fine-tuned promptable baselines: SAM-Med2D \cite{b2}, MedSAM \cite{b9}, and ScribblePrompt-UNet \cite{b11}. Attention U-Net is trained without a prompt, to show how hard the task becomes when localization must be fully automatic. Among the promptable baselines, SAM-Med2D is fine-tuned and evaluated at $256\times256$ using the same dataset splits and the same box prompts as PGA-UNet, so PGA-UNet at $256$ against SAM-Med2D at $256$ is the article's resolution-matched prompt comparison. MedSAM and ScribblePrompt-UNet run at their own native resolutions ($1024$ and $128$); Table~\ref{tab:sam-ref} reports these alongside PGA-UNet's own $128$ and $512$ results as configuration-specific reference results, each scored in its own frame, and does not use them to establish a cross-model ranking.

To compare the proposed design with conventional ways of using the same box, we evaluate two additional baselines based on the Attention U-Net backbone, without PGA-UNet's Gaussian-prior heatmap, Prompt Spatial Gate, or Conditional Attention Decoder. The channel baseline receives a binary box mask as an additional input channel, whereas the crop baseline receives the image region defined by the box. Only the crop baseline retains the backbone's original attention-gated skip connections unmodified, since the box only changes what image region enters an otherwise-standard Attention U-Net, so we refer to it as Attention U-Net + prompt crop; the channel baseline instead concatenates the mask at the input and decodes through plain, non-gated skip connections, so it is not attention-gated in the same sense as the automatic Attention U-Net reference, and we name it accordingly, U-Net + binary prompt channel. The full PGA-UNet model uses a Gaussian-smoothed heatmap together with explicit encoder and decoder conditioning. These comparisons therefore concern the complete tested configurations, combining differences in prompt representation, backbone gating, and PGA-UNet's own conditioning mechanism, not any one of these in isolation. Both baselines are trained under the same covering and off-center box protocol as PGA-UNet. The binary-prompt PGA-UNet ablation (Section~\ref{sec:ablation}) provides a comparison using the same binary box representation as the channel baseline, without isolating every architectural component. When comparing models against each other, we use the pasted-back full-frame metric for the prompt-crop baseline, not the diagnostic value measured inside the crop frame.

Every comparison table in this article reports Dice and IoU for region overlap, CBL for localization, and HD95 for boundary distance. Dice and IoU are fairly strongly correlated overlap measures; we place them side by side only for the reader's convenience in cross-checking. On the small-lesion subsets specifically, IoU has a low absolute value and is quite sensitive to a few boundary pixels, so Dice is the primary overlap metric used there. Let $(x_p,y_p)$ be the predicted mask centroid, $(x_g,y_g)$ the ground-truth centroid, and $D_{\mathrm{GT}}$ the diagonal length of the ground-truth bounding box; CBL is defined as follows:
\begin{equation}
\mathrm{CBL}=\max\left(0,1-\frac{\sqrt{(x_p-x_g)^2+(y_p-y_g)^2}}{D_{\mathrm{GT}}}\right).
\end{equation}
HD95 is computed from bidirectional boundary distances. When both masks are empty, HD95 is set to 0; if only one mask is empty, HD95 is set to the input side length $S$ (128, 256, or 512 pixels, depending on the configuration). No sample is discarded from the computation. For comparisons spanning resolutions, we additionally report HD95 in normalized form, $\mathrm{HD95}/(\sqrt{2}\,S)$, since a raw pixel distance cannot be directly compared across the $128$, $256$, and $512$ frames. Some subgroup comparisons mix models trained at $256\times256$ and $512\times512$; in that case, everything is scored together in a shared $512\times512$ frame, with the $256\times256$ predictions upsampled before scoring, so every model is measured against exactly the same ground-truth mask, rather than in its own native frame. Cross-model comparisons used to support the main conclusions are restricted to a shared scoring frame: the conventional prompt baselines and PGA-UNet at $512\times512$ (Table~\ref{tab:robust}), and SAM-Med2D and PGA-UNet at $256\times256$ (Table~\ref{tab:sam}). The lower-area subset comparison with SAM-Med2D (Table~\ref{tab:small-sam}) uses $256\times256$ model inputs and a shared $512\times512$ scoring frame. Other operating configurations, including PGA-UNet's own $128$ and $512$ results and the MedSAM and ScribblePrompt-UNet results, are documented separately in Table~\ref{tab:sam-ref} as configuration-specific reference results, each scored in its own frame (PGA-UNet at $128$ or $512$; MedSAM and ScribblePrompt-UNet at the original image resolution, not at the $128$ or $1024$ frame their own network internally resizes to); their raw HD95 values are shown for completeness but are left unbolded. Dice, IoU, and CBL are dimensionless, but they need not remain unchanged after mask resampling and rasterization: a model run at a lower native resolution discretizes a lesion's boundary into fewer pixels before any of these metrics is computed. Consequently, results scored in different frames are not used to establish a cross-model ranking in this article. Normalizing HD95 changes its distance scale but does not remove differences introduced by prediction and ground-truth resampling.

For the training-free self-assessment score, each lesion prompt yields a score $Q$ together with the true thresholded Dice of its own predicted mask; we report the Pearson and Spearman correlation between the two, separately for the covering and off-center conditions. For the candidate-region shortlist, each test image is sampled with 50 boxes whose side length is uniform in $[0.15,0.55]$ of the frame, each box is scored by $Q$, and the five highest-scoring boxes are kept; each retained box is then scored by coverage, the fraction of ground-truth lesion pixels it contains, and no predicted mask enters this coverage measure. From the same 50-box pool we also compute a random-5 baseline (mean over 2,000 resamples per image, without replacement) and an oracle-5 upper bound (the 5 boxes with the truly highest coverage), to check whether $Q$-ranking adds value over chance. As a negative control, the same 50-box sampling is also run on the lesion-free images of each dataset (1879 for BTXRD, 3366 for FracAtlas); with no lesion present there is no coverage to compute, so we report only the distribution of $Q$.

\subsection{Statistical Analysis}
The main baseline comparisons (Tables~\ref{tab:baseline}, \ref{tab:sam}, \ref{tab:extreme}) report descriptive point estimates on the fixed test splits; no paired uncertainty estimates are reported for these comparisons. The existing ablation analysis (Section~\ref{sec:ablation}) uses paired per-image Dice values with a two-sided sign-flip randomization test and percentile bootstrap confidence intervals, each based on $20{,}000$ resamples, with Holm correction as specified in Table~\ref{tab:ablation-sig}. These results are conditional on the evaluated split and trained checkpoints. Four additional image-level splits (Section~\ref{sec:mccv}) characterize variation in PGA-UNet's own performance while keeping its training seed fixed; they do not estimate variation due to the training seed itself, nor do they establish the stability of differences from competing models. The exploratory self-assessment correlation analysis (Section~\ref{sec:selfassess-results}) retains its own permutation $p$-values and bootstrap $95\%$ intervals for Spearman's $\rho$. All resampling uses the fixed seed 22120196.

\subsection{Scope of Analysis}
This article focuses only on prompt-guided segmentation itself. A lesion-screening classification stage was explored in a separate study and is not part of this article. The self-assessment score and the candidate-region shortlist are treated only as auxiliary analyses, within the same narrow scope already stated in Section~\ref{sec:selfassess}; apart from one negative control on lesion-free images, they are not extended any further here. To keep the article compact, we retain only the subgroup analyses that most directly support the main claim: the role of prompt-guided localization (Section~\ref{sec:extreme}), sensitivity to controlled prompt displacement as defined above (Section~\ref{sec:robust}), and observed results on the images with the lowest total lesion area (Section~\ref{sec:small}, where the subgroup definition and its coverage of each test split are given). It should also be reiterated: partial-coverage boxes, negative boxes, and off-target boxes are outside the scope of this study, so the article's claims do not rely on them.
